\documentclass[10pt,a4paper]{amsart}
\usepackage[margin=19mm]{geometry}
\usepackage{amsmath,amssymb,mathtools}
\usepackage[scr=rsfs,cal=euler]{mathalfa}
\usepackage{xfrac}
\usepackage[unicode,hidelinks,bookmarks=false]{hyperref}
\newcommand{\ie}{i.e.\ }
\newcommand{\cf}{cf.\ }
\newcommand{\resp}{respectively\ }
\newcommand{\Subsection}{\subsection}
\providecommand{\nobreakdash}{\nobreak\hbox{-}}
\setlength{\emergencystretch}{2em}
\multlinegap0pt
\usepackage{amssymb}
\usepackage{tikz}
\newtheorem{lem}{Lemma}
\newtheorem{pf}{Pf}
\title{The Bohr Phenomenon for analytic functions on simply connected domains}
\author{Molla Basir Ahamed, Vasudevarao Allu,, Himadri Halder}
\date{Source: arXiv:2011.13890v1 (CC BY 4.0)}
\begin{document}
\maketitle

\noindent\emph{Source and licence.} The Bohr Phenomenon for analytic functions on simply connected domains. Version arXiv:2011.13890v1 (CC BY 4.0), distributed by arXiv under the Creative Commons Attribution 4.0 International licence, http://creativecommons.org/licenses/by/4.0/, as recorded in the arXiv licence metadata for this exact version and shown on the abstract page https://arxiv.org/abs/2011.13890v1. Full licence notice: \texttt{licenses/arxiv-2011.13890-CC-BY-4.0.txt}.\par\smallskip

\noindent\textit{Editorial scope.} Counted unit: the article's lem lem-2.7 (source0.tex lines 369--375). The article's complete original proof of it is delivered with it (source0.tex lines 527--587, one \texttt{proof} environment of 61 lines, which the article places away from the statement). No mathematics is written, repaired or re-derived: the statement and the proof are the article's own lines, in the article's own notation, and no retained line is substituted or re-flowed.  No further source-local context is needed: every cross-reference of every retained line resolves inside the two delivered spans.  Nothing was omitted from any retained span, so the package carries no omission record.  No retained line contains a citation, so the wrapper carries no bibliography and no bibliography entry.  No retained line contains a figure environment, an image-inclusion command or drawing code, so no image is delivered and the package declares no asset dependency.  Editorial formatting only, in addition: an AMS \texttt{amsart} preamble that restates the article's own declarations and macros used by the retained lines, namely the environments \texttt{lem}, \texttt{pf} on separate counters, together with the standard packages those lines need.  The retained environments print in this document as Lemma 1, Pf 1 in order of appearance, whereas the article prints the same items with the numbering of its own full document.  The complete native source of the article as distributed in the arXiv e-print for this exact version is bundled unchanged as \texttt{sources/105164/source0.tex}.
\begin{lem}\label{lem-2.7}
	Let $ g : \mathbb{D}\rightarrow\overline{\mathbb{D}} $ be an analytic function, $ \lambda\in [0,512/243] $ and let $ \gamma\in\mathbb{D} $ be such that $ g(z)=\sum_{n=0}^{\infty}\alpha_n(z-\gamma)^n $ for $ |z-\gamma|<1-|\gamma| $. Then 
	\begin{align*}
		\sum_{n=0}^{\infty}|\alpha_n|\rho^n+\left(\frac{8}{9}-\frac{27}{64}\lambda\right)\left(\frac{S^{\gamma}_{\rho}}{\pi}\right)+\lambda\left(\frac{S^{\gamma}_{\rho}}{\pi}\right)^2\leq 1\;\; \text{for}\;\; \rho\leq \rho_0=\frac{1-|\gamma|^2}{3+|\gamma|},
	\end{align*}
	where $ S^{\gamma}_{\rho} $ denotes the area of the image of the disk $ \mathbb{D}(\gamma;r(1-|\gamma|)) $ under the mapping  $ g $.
\end{lem}

\begin{pf}[\bf Proof of Lemma \ref{lem-2.7}]
	Without loss of generality, we assume that $ \gamma\in [0,1) $. Also let $ z\in\mathbb{D}_{\gamma}:=\mathbb{D}(\gamma;1-\gamma) $ if, and only if, $ w=(z-\gamma)/(1-\gamma)\in\mathbb{D}. $
	Then we have 
	\begin{align*}
		g(z)=\sum_{n=0}^{\infty}\alpha_n(1-\gamma)^n\phi^n(z)=\sum_{n=0}^{\infty}b_n\phi^n(z):=G(\phi(z))
	\end{align*}
	for $ z\in \mathbb{D}_{\gamma}, $ where $ b_n= \alpha_{n} (1-\gamma)^n$. A simple computation shows that
	\begin{align}\label{e-2.8}
		\frac{S^{\gamma}_{\rho}}{\pi}=\frac{1}{\pi}\text{Area}\bigg(G(\mathbb{D}(0,\rho))\bigg)\leq (1-|b_0|^2)^2\frac{\rho^2}{(1-\rho^2)^2}=(1-|\alpha_0|^2)^2\frac{\rho^2}{(1-\rho^2)^2}.
	\end{align}
	Therefore,
	\begin{align}\label{e-2.9}
		\sum_{n=1}^{\infty}|\alpha_n|\rho^n\leq\frac{1-|\alpha_0|^2}{1+\gamma}\sum_{n=1}^{\infty}\left(\frac{\rho}{1-\gamma}\right)^n=\frac{1-|\alpha_0|^2}{1+\gamma}\frac{\rho}{1-\gamma-\rho}.
	\end{align}
	In view of \eqref{e-2.8} and \eqref{e-2.9}, we obtain 
	\begin{align*} &
		|\alpha_0|+	\sum_{n=1}^{\infty}|\alpha_n|\rho^n+k\left(\frac{S^{\gamma}_{\rho}}{\pi}\right)+\lambda\left(\frac{S^{\gamma}_{\rho}}{\pi}\right)^2\\&= |\alpha_0|+\frac{(1-|\alpha_0|^2)\rho}{(1+\gamma)(1-\gamma-\rho)}+k\frac{(1-|\alpha_0|^2)^2\rho^2}{(1-\rho^2)^2}+\lambda\frac{(1-|\alpha_0|^2)^4\rho^4}{(1-\rho^2)^4}\\&= 1+\Psi_1^{\gamma}(\rho),
	\end{align*}
	where 
	\begin{align*}
		\Psi_1^{\gamma}(\rho)=\frac{(1-|\alpha_0|^2)\rho}{(1+\gamma)(1-\gamma-\rho)}+k\frac{(1-|\alpha_0|^2)^2\rho^2}{(1-\rho^2)^2}+\lambda\frac{(1-|\alpha_0|^2)^4\rho^4}{(1-\rho^2)^4}-(1-|\alpha_0|)
	\end{align*}
	which can be written as 
	\begin{align*} 
		\Psi_1^{\gamma}(\rho)&=\frac{1-|\alpha_0|^2}{2}\bigg(1+2\lambda(1-|\alpha_0|^2)^3 \left(\frac{\rho^4}{(1-\rho^2)^4}+\frac{k}{\lambda}\frac{\rho^2}{(1-\rho^2)^2(1-|\alpha_0|^2)^2}\right)\\&\;\;\;\;+\frac{1}{2\lambda(1-|\alpha_0|^2)^3}\left(\frac{2\rho}{(1+\gamma)(1-\gamma-\rho)}-1\right)-\frac{2}{1+|\alpha_0|}\bigg).
	\end{align*}
	Let $ \rho\leq \rho_0 $. Then, it is easy to see that $\Psi_1^{\gamma}(\rho)$ is an increasing function and hence $ 	\Psi_1^{\gamma}(\rho)\leq 	\Psi_1^{\gamma}(\rho_0) $, where 
	\begin{align*}
		\frac{2\rho_0}{(1+\gamma)(1-\gamma-\rho_0)}=1,\;\;\text{\it i.e.,}\;\; \rho_0=\frac{1-\gamma^2}{3+\gamma}.
	\end{align*}
	A simple computation shows that 
	\begin{align*}
		\Psi_1^{\gamma}(\rho_0)&=\frac{1-|\alpha_0|^2}{2}\left(1+2\lambda(1-|\alpha_0|^2)^3A^4(\gamma)+2k(1-|\alpha_0|^2)A^2(\gamma)-\frac{2}{1+|\alpha_0|}\right)\\&=\frac{1-|\alpha_0|^2}{2}J(|\alpha_0|),
	\end{align*}
	where 
	\begin{align*}
		J(x)&=1+2\lambda(1-x^2)^3A^4(\gamma)+2k(1-x^2)A^2(\gamma)-\frac{2}{1+x}\;\; \mbox{for}\;\ x\in [0,1]\\\text{and}\;\; A(\gamma)&=\frac{(3+\gamma)(1-\gamma^2)}{(3+\gamma)^2-(1-\gamma^2)^2}.
	\end{align*}
	It is enough to show that $ J(x)\leq 0 $ for $ x\in [0,1] $ and $ \gamma\in [0,1) $ so that $ 	\Psi_1^{\gamma}(\rho_0)\leq 0 $. We note that $ A(\gamma)>0 $ for $ \gamma\in [0,1) $. Further,
	\begin{align*}
		J(0)=2\lambda A^4(\gamma)+2kA^2(\gamma)-1, \;\; \text{and}\;\; \lim_{x\rightarrow 1^{-}}J(x)=0.
	\end{align*}
	It can be seen that $ A(\gamma)=(f_1\circ f_2)(\gamma) $, where $ f_1(\rho)=\rho/(1-\rho^2) $ and $ f_2(\gamma)=(1-\gamma^2)/(3+\gamma). $
	Since $ A^{\prime}(\gamma)=f^{\prime}_1(f_2(\gamma)f^{\prime}_2(\gamma)), $ where 
	\begin{equation}
		f^{\prime}_2(\gamma)=-\left(\frac{\gamma^2+6\gamma+1}{(3+\gamma)^2}\right)<0
	\end{equation}
	which implies that $ f_1(\rho) $ is an increasing function of $ \rho $ in $ (0,1) $, and $ f_2 $ is a decreasing function of $ \gamma $ in $ [0,1) $. Hence, it follows that $ A(\gamma) $ is a decreasing function of $ \gamma $  in $ [0,1) $, with $ A(0)=3/8 $ and $ A(1)=0 $. It can be seen that $ A^2(\gamma) $ and $ A^4(\gamma) $ are decreasing functions on $ [0,1) $. Therefore, we have 
	\begin{align*}
		A^2(\gamma)\leq A^2(0)=\frac{9}{64}\;\; \text{and}\;\; A^4(\gamma)\leq A^4(0)=\frac{81}{4096}.
	\end{align*}
	Since $ x\in [0,1] $, we have
	\begin{align*}
		x(1+x)^2A^2(\gamma)\leq \frac{9}{16}\;\;\text{and}\;\; x(1+x)^2(1-x^2)^2A^4(\gamma)\leq \frac{81}{1024}.
	\end{align*}
	As a consequence, we obtain
	\begin{align*}
		J^{\prime}(x)&=\frac{2}{(1+x)^2}\bigg(1-2kx(1+x)^2A^2(\gamma)-6\lambda x(1+x)^2(1-x^2)^2A^4(\gamma)\bigg)\\&\geq  \frac{2}{(1+x)^2}\bigg(1-\left(\frac{9k}{8}+\frac{243\lambda}{512}\right)\bigg)\\&\geq 0,\;\;\;\;\; \text{if}\;\; k+27\lambda/64\leq 8/9.
	\end{align*}
	Therefore, $ J(x) $ is an increasing function in $ [0,1] $ for $ k+27\lambda/64\leq 8/9 $. Hence, $ J(x)\leq 0 $ for all $ x\in [0,1] $ and $ \gamma\in [0,1) $. This completes the proof.
\end{pf}

\end{document}
